% Put your abstract here

\begin{abstractwithpageno}
International School Manila has four separate systems for handling sponsorship billing and Letter of Guarantee (LoG) coverage: PowerSchool (SIS), the Student Charging Portal (SCP), NetSuite (ERP), and the Online Billing System (OBS). Prior to the development of this project, coverage decisions were manually interpreted at charge entry and this could lead to inconsistent assignments to \textbf{Bill-To}, causes for delay in the correction process, lack of audit trails and further reconciliation efforts. The project created and implemented a pilot \textbf{Sponsor Management and LoG Coverage Integration System} that combines all the master data of sponsors, tracks LoG records, handles workflows based on the roles and assesses coverage using the \textbf{deterministic business rules}.

The \textbf{deployed pilot web application} was done with ASP.NET Core including REST API endpoints, role-based authentication for Administrator, Admissions and Cashier users, local credentials for authentication for ISM staff accounts, and Azure deployment for production-style validation. The system enables the management of sponsor profiles, prevention of sponsor duplication, creation and activation of LoG, sponsor change requests, sponsor views, coverage preview, creation of audit trail and reporting support. The web interface has been responsive and is \textit{PWA-ready}, but has not been fully developed as a Progressive Web App as the \textit{PWA installable features}, \textit{service worker caching}, \textit{offline behavior} and the \textit{web app manifest} were not in the scope of the pilot.

A total of \textbf{110 test cases} were found in total, of which \textbf{95 passed}, \textbf{12 noted} and \textbf{3 known non-critical items} were noted; none of them were critical or blocking. \textbf{60 user acceptance testing (UAT) workflow scenarios} were verified; \textbf{54 were successful} and \textbf{6 were non-blocking issues} for refinement. The proven workflows covered sponsor search, sponsor maintenance, creation and viewing of the LoG, coverage check, restricted access, report generation, and sponsor self-service actions. There are minor usability and validation issues such as the delayed LoG reassignment display after sponsor merge, the lack of visibility for attachments in one change-request flow, unclear confirmation messages, incomplete report header metadata, and improved sponsor access-error messages.

The pilot proved to show that existing manual interpretation can be minimised, governance is enhanced, audit preparedness is improved and a practical roadmap towards production integration with current ISM systems is achieved with a centralized \textbf{Sponsor Master} and \textbf{deterministic coverage engine}. The pilot web application deployed is a working example that shows that there is no need to replace ISM's existing core platforms, and that consistent and explainable decisions on coverage of sponsors can be made. This pilot web application is a working example that demonstrates that consistent and explainable decisions on coverage of sponsors can be made without replacement of existing ISM core platforms.
\end{abstractwithpageno}